% QUESTAO
A área do círculo e o comprimento da circunferência no caso abaixo são:

\parbox{0.6\linewidth}{

% ALTERNATIVAS
$36\pi$\,m$^2$ e $12\pi$\,m
$25\pi$\,m$^2$ e $10\pi$\,m
$20\pi$\,m$^2$ e $5\pi$\,m
$30\pi$\,m$^2$ e $20\pi$\,m
$25\pi$\,m$^2$ e $15\pi$\,m

}\parbox{0.39\linewidth}{
\begin{center}
\vspace{-9mm}
	\begin{tikzpicture}[>=latex,scale=1.1,rotate=0]
	\coordinate (O) at (0,0);   \fill (O) circle (0.4mm);
	\coordinate (A) at (80:1);
	\coordinate (B) at (-100:1);
	\draw[thick] (O) circle (1);
	\draw[thick] (A)--(B);
	
	\draw[ decorate, decoration={brace,raise=3pt} ]  (80:0.95) -- (-100:0.95);
	\draw[nodelabel={18pt}{$12$\,m}] (A) -- (B);
	
	\end{tikzpicture}
\vspace{-9mm}
\end{center}
}


